\chapter{Verificatie, simulatie en oplevering}
\label{chap:testen}

\section{Teststrategie}
Een groen testresultaat is alleen waardevol wanneer de verwachte uitkomst onafhankelijk van de implementatie is bepaald. Daarom gebruiken de tests expliciete scenario's met handmatig berekende tijden, grenzen en negatieve gevallen. Ze controleren niet enkel dat een functie een waarde retourneert, maar ook welke ronden worden geselecteerd, welke eenheid wordt gebruikt en wanneer een resultaat bewust onbekend moet blijven.

De testsuite wordt uitgevoerd met Node.js en \code{assert}. \code{tests/run-all.js} ontdekt automatisch alle bestanden met suffix \code{.test.js}. Daardoor volstaat een nieuw testbestand om het in de volledige suite op te nemen. De suite bevat momenteel tests voor parser, storage schema, sessieherstel, analytics, voorspelling, referentietijden, klassegevechten, pitstopplanning, circuitprofielen, grafiekdata, sessiemodi, rendererviewmodellen, updater en applicatielifecycle.

\section{Mockdata en historische race}
Voor statistische functies werd een centrale mockhistoriek gemaakt met meerdere piloten, stints, sectoren en ongeldige ronden. Hiermee worden gemiddelde rondetijden, snelste ronden, sectorgemiddelden, beste sectoren en vergelijkingen met klassewagens gecontroleerd. Edgecases omvatten lege reeksen, nullwaarden, onbekende piloten, negatieve of onrealistische tijden, gedeeltelijke sectorinformatie en wissels tussen piloten.

Daarnaast werd historische data van een Belcar-race in Assen gebruikt. PDF-bestanden met startgrid, klassering, rondetijden en sectoranalyse werden omgezet naar een compacte dataset voor de Club Challenge-klasse. Een lokale HTTP-simulator presenteert deze data als een timingpagina en kan de race versneld laten verlopen, bijvoorbeeld één ronde per vijf of tien seconden. De simulator heeft eigen tests; afzonderlijke analytics-tests gebruiken dezelfde historische data om de berekeningsfuncties te controleren.

Deze twee soorten tests hebben een ander doel. Simulatortests bewijzen dat de nepwebsite correct evolueert. Analytics-tests bewijzen dat domeinfuncties de verwachte uitkomst geven. De simulator alleen kan geen fout in een berekening detecteren wanneer dezelfde fout ook in de verwachte simulatorlogica zou zitten.

\section{Gevonden fouten door realistische scenario's}
Een aantal belangrijke problemen werd alleen zichtbaar met realistische data:

\begin{itemize}
  \item Bij een nieuwe sessiemap verscheen aanvankelijk ronde 60, hoewel de live demo verder stond. De oorzaak was dat wijzigingen in de laatste rondetijd als nieuwe lokale rondes werden geteld zonder de echte rondecontext correct te gebruiken.
  \item Een prediction van ongeveer 1:35 voor een wagen met ronden rond 2:05 wees op fout geïnterpreteerde sectorwaarden en onvoldoende samplevalidatie.
  \item Een klassewagen leek slechts twaalf seconden achter te liggen, terwijl een tussenliggende wagen één ronde achterstand had. De oplossing gebruikt de algemene rangschikking en alle tussenliggende intervallen.
  \item Een pitstopprojectie meldde duizenden seconden voorsprong doordat een rondeveld als tijd werd geïnterpreteerd. Provider- en eenheidsdetectie werden aangescherpt.
  \item Safety Car- en FCY-ronden maakten grafieken onleesbaar en gemiddelden fout. De gedeelde eligibilityfuncties sluiten ze nu overal consistent uit.
  \item Een verborgen livevenster blokkeerde op macOS het normaal afsluiten van de app. Een geteste lifecyclemodule maakt bij Quit alle verborgen vensters echt afsluitbaar.
\end{itemize}

Deze voorbeelden tonen waarom alleen happy-path-tests onvoldoende zijn. Veel fouten ontstaan precies wanneer een veld syntactisch geldig is maar semantisch iets anders betekent.

\section{Coverage en continue integratie}
Naast \code{npm test} is een coveragerun beschikbaar via \code{npm run test:coverage}. Statement-, function- en branchcoverage worden gebruikt als aanwijzing voor vergeten paden. Vooral branchcoverage is nuttig voor fallbacklogica, nullwaarden, providerkeuzes en gesloten/open pitvensters. Een hoog percentage bewijst echter niet dat de verwachtingen correct zijn; daarom blijft scenarioselectie belangrijker dan het getal alleen.

GitHub Actions voert bij elke push en pull request zowel tests als coverage uit op Ubuntu met Node.js. Branch protection kan deze workflow als verplichte statuscheck gebruiken voordat een pull request naar main wordt gemerged.

\section{Distributie en updates}
\code{electron-builder} maakt een NSIS-installer voor Windows en DMG/ZIP-bestanden voor macOS. Een tag zoals \code{v1.1.1} activeert de releaseworkflow. Mac- en Windows-builds worden parallel gemaakt, waarna één GitHub Release met alle platformbestanden wordt gepubliceerd.

De geïnstalleerde applicatie controleert bij het starten via \code{electron-updater} op een nieuwere GitHub Release. Developmentruns via \code{npm start} doen dit bewust niet. De app en installers zijn momenteel niet digitaal ondertekend of genotariseerd. Windows SmartScreen en macOS Gatekeeper kunnen daarom waarschuwingen tonen. De workflow bevat reeds plaatsen voor toekomstige signingsecrets.

\section{Beperkingen van de huidige verificatie}
Automatische tests modelleren bekende scenario's. Ze vervangen geen langdurige test met echte timingproviders, netwerkonderbrekingen en de operationele gebruiker. De Assen-data is waardevol, maar vertegenwoordigt niet alle circuits, reglementen en timingimplementaties. De geplande test in Spa moet daarom de volledige keten valideren: installatie, polling, opslag, analyses, pitstopprojectie, meerdere schermen en interpretatie door de opdrachtgever.

